% !TEX root = ../main.tex
%% Chapitre 5 — Logarithme népérien
%% Sources : frontend/cours/resume/terminal/resume_fonction_logarithme_terminal.txt
%%           frontend/exercices/terminal_spé_Maths/exercice_fonction_logarithme.txt
%% Repris des chapitres 2 à 4 : questions avec ln et étude de ln((x²+1)/(x²-4)).
%% Exercices 12 à 14 : problèmes inédits rédigés pour le livre.

\begin{ouverture}[ln]{5}{Logarithme népérien}{Partie 1 · Analyse}
  {Définition\enspace·\enspace Propriétés algébriques\enspace·\enspace Dérivée\enspace·\enspace
   Limites\enspace·\enspace Équations}
  \objectif{Utiliser les propriétés algébriques du logarithme pour transformer une expression.}
  \objectif{Résoudre des équations et des inéquations, en tenant compte de l'ensemble de définition.}
  \objectif{Résoudre une inéquation du type $q^{n}<c$ pour déterminer un seuil.}
  \objectif{Dériver une fonction de la forme $\ln(u)$ et étudier une fonction avec $\ln$.}
  \objectif{Utiliser les croissances comparées du logarithme et des puissances.}
  \prerequis{Fonction exponentielle : propriétés et limites (Première, chapitre 2).}
  \prerequis{Fonctions composées et dérivation (chapitre 3).}
  \prerequis{Continuité et théorème des valeurs intermédiaires (chapitre 4).}
\end{ouverture}

%% =====================================================================
%%  L'ESSENTIEL
%% =====================================================================
\essentiel

\essSection{1}{Définition}

\begin{definition}[Logarithme népérien]
Pour tout réel $x>0$, l'équation $e^{y}=x$ admet une unique solution réelle $y$, notée $\ln x$. La fonction
$\ln$, définie sur $\intOO{0}{+\infty}$, est la \textbf{fonction logarithme népérien}. Ainsi :
\[ e^{\ln x}=x\ \text{ pour tout } x>0 \qquad\qquad \ln\left(e^{x}\right)=x\ \text{ pour tout réel } x \qquad\qquad \ln1=0,\enspace \ln e=1. \]
\end{definition}

\noindent
\begin{minipage}[c]{0.5\linewidth}
Les fonctions $\exp$ et $\ln$ sont \textbf{réciproques} l'une de l'autre : dans un repère orthonormé, leurs courbes
sont symétriques par rapport à la droite d'équation $y=x$. Le point $(0\pv1)$ de la courbe de $\exp$ correspond au
point $(1\pv0)$ de celle de $\ln$, et $(1\pv e)$ correspond à $(e\pv1)$.
\par\vspace{3mm}
\begin{savoirfaire}[Équations et inéquations]
Pour $a>0$ et $b>0$ : $\ln a=\ln b\iff a=b$ et $\ln a<\ln b\iff a<b$.\\
Pour tout réel $k$ : $\ln x=k\iff x=e^{k}$.\\
\emph{On commence toujours par déterminer l'ensemble de résolution.}
\end{savoirfaire}
\end{minipage}\hfill
\begin{minipage}[c]{0.45\linewidth}
\centering
\begin{tikzpicture}
\begin{axis}[otGraphe, width=70mm, xmin=-2.6, xmax=4.4, ymin=-2.6, ymax=4.4,
  xtick={1,2,3}, ytick={1,2,3}]
  \addplot[otMuted,dashed,line width=0.5pt,domain=-2.6:4.3] {x};
  \addplot[otNavy,line width=1.1pt,domain=-2.6:1.45,samples=80] {exp(x)};
  \addplot[otGreen,line width=1.2pt,domain=0.075:4.3,samples=120] {ln(x)};
  \addplot[only marks,mark=*,mark size=1.4pt,otNavy] coordinates {(0,1) (1,2.718)};
  \addplot[only marks,mark=*,mark size=1.4pt,otGreen] coordinates {(1,0) (2.718,1)};
  \node[anchor=east,font=\scriptsize,text=otNavy] at (axis cs:1.3,3.9) {$\exp$};
  \node[anchor=north,font=\scriptsize,text=otGreenDark] at (axis cs:3.9,1.2) {$\ln$};
  \node[anchor=south west,font=\scriptsize,text=otMuted] at (axis cs:3.2,3.5) {$y=x$};
\end{axis}
\end{tikzpicture}\\[0.5mm]
\legende{Courbes de $\exp$ et de $\ln$, symétriques par rapport à $y=x$.}
\end{minipage}

\essSection{2}{Propriétés algébriques}

{\renewcommand{\arraystretch}{1.8}\arrayrulecolor{otRule}%
\begin{tabularx}{\linewidth}{@{}>{\otSemi\fontsize{7.5}{10}\selectfont\leavevmode\color{otNavy}}p{30mm}YY@{}}
\toprule
 & {\otSemi\fontsize{8.5}{11}\selectfont Pour tous réels $a>0$ et $b>0$} & {\otSemi\fontsize{8.5}{11}\selectfont Exemple} \\
\midrule
\otCaps[8]{Produit} & $\ln(ab)=\ln a+\ln b$ & $\ln6=\ln2+\ln3$ \\
\otCaps[8]{Quotient} & $\ln\left(\dfrac ab\right)=\ln a-\ln b$ & $\ln\left(\dfrac12\right)=-\ln2$ \\
\otCaps[8]{Puissance} & $\ln\left(a^{n}\right)=n\ln a$\enspace($n$ entier relatif) & $\ln8=3\ln2$ \\
\otCaps[8]{Racine carrée} & $\ln\sqrt a=\dfrac12\ln a$ & $\ln\sqrt e=\dfrac12$ \\
\bottomrule
\end{tabularx}}

\begin{demonstration}[de la relation fonctionnelle]
$e^{\ln a+\ln b}=e^{\ln a}\times e^{\ln b}=ab=e^{\ln(ab)}$. La fonction exponentielle étant strictement croissante,
deux réels qui ont la même exponentielle sont égaux : $\ln(ab)=\ln a+\ln b$.
\end{demonstration}

\essSection{3}{Étude de la fonction logarithme}

\begin{theoreme}[Dérivée et variations]
La fonction $\ln$ est dérivable sur $\intOO{0}{+\infty}$ et $\ln'(x)=\dfrac1x$. Elle est donc strictement croissante sur
$\intOO{0}{+\infty}$, avec $\displaystyle\lim_{x\to0^{+}}\ln x=-\infty$ et $\displaystyle\lim_{x\to+\infty}\ln x=+\infty$.
\end{theoreme}

\noindent
\begin{minipage}[c]{0.5\linewidth}
\centering
\begin{tableau}
  \otTab[1.5]{$x$/0.9,$\ln'(x)$/0.9,$\ln$/1.9}{$0$,$1$,$e$,$+\infty$}
  \tkzTabLine{d,+,t,+,t,+,}
  \tkzTabVar{D-/ /$-\infty$,R/,R/,+/$+\infty$}
  \tkzTabIma{1}{4}{2}{$0$}
  \tkzTabIma{1}{4}{3}{$1$}
\end{tableau}
\par\vspace{3mm}
\raggedright
La courbe est sous sa tangente en $1$, d'équation $y=x-1$ ($\ln$ est concave car $\ln''(x)=-\frac{1}{x^{2}}<0$) :
\[ \ln x\leq x-1\ \text{ pour tout } x>0. \]
\end{minipage}\hfill
\begin{minipage}[c]{0.45\linewidth}
\centering
\begin{tikzpicture}
\begin{axis}[otGraphe, width=70mm, xmin=-0.6, xmax=5.4, ymin=-3.2, ymax=2.4,
  xtick={1,2,3,4,5}, ytick={-2,-1,1,2}]
  \addplot[otMuted,dashed,line width=0.6pt,domain=-0.3:3.3] {x-1};
  \addplot[otNavy,line width=1.2pt,domain=0.042:5.3,samples=150] {ln(x)};
  \addplot[only marks,mark=*,mark size=1.4pt,otNavy] coordinates {(1,0) (2.718,1)};
  \node[anchor=north west,font=\scriptsize] at (axis cs:1,0) {$1$};
  \node[anchor=south east,otEtiquette,font=\scriptsize] at (axis cs:2.7,1.05) {$(e\pv1)$};
  \node[anchor=west,font=\scriptsize,text=otMuted] at (axis cs:2.6,2.15) {$y=x-1$};
\end{axis}
\end{tikzpicture}\\[0.5mm]
\legende{L'axe des ordonnées est asymptote verticale.}
\end{minipage}

\begin{theoreme}[Logarithme d'une fonction]
Si $u$ est dérivable et strictement positive sur $I$, alors $\ln(u)$ est dérivable sur $I$ et
$\big(\ln(u)\big)'=\dfrac{u'}{u}$.
\end{theoreme}

\essSection{4}{Limites et croissances comparées}

\begin{theoreme}[Limites de référence]
\[
  \lim_{x\to+\infty}\frac{\ln x}{x}=0 \qquad\quad
  \lim_{x\to0^{+}}x\ln x=0 \qquad\quad
  \lim_{x\to+\infty}\frac{\ln x}{x^{n}}=0\ (n\geq1) \qquad\quad
  \lim_{x\to0}\frac{\ln(1+x)}{x}=1
\]
\end{theoreme}

En $+\infty$, toute puissance de $x$ « l'emporte » sur le logarithme. La dernière limite est le nombre dérivé de
$x\mapsto\ln(1+x)$ en $0$.

\begin{savoirfaire}[Déterminer un seuil avec le logarithme]
Pour $0<q<1$ et $c>0$ : $q^{n}<c\iff n\ln q<\ln c\iff n>\dfrac{\ln c}{\ln q}$.\enspace
On divise par $\ln q<0$ : \textbf{le sens de l'inégalité change}.
\end{savoirfaire}

\begin{attention}[Pièges fréquents]
\begin{itemize}[label={\color{otOrange}\textbullet},leftmargin=1.4em,itemsep=0.6mm,topsep=0.5mm]
  \item $\ln(a+b)\neq\ln a+\ln b$ : c'est $\ln(ab)$ qui vaut $\ln a+\ln b$.
  \item $\ln x$ n'existe que pour $x>0$ : vérifier les solutions trouvées par rapport à l'ensemble de définition.
  \item $\ln q<0$ lorsque $0<q<1$ : diviser par $\ln q$ change le sens d'une inégalité.
  \item $\big(\ln u\big)'=\dfrac{u'}{u}$, et non $\dfrac1u$.
\end{itemize}
\end{attention}

%% =====================================================================
%%  EXERCICES
%% =====================================================================
\exercices
\rubrique{Découverte}

\begin{exercice}[1]{Propriétés algébriques}
Écrire chacune des expressions suivantes sous la forme d'un seul logarithme, ou la simplifier.
\begin{tasks}(3)
  \task $\ln8+\ln2$
  \task $\ln27-\ln3$
  \task $2\ln5+3\ln2$
  \task $\ln\left(\dfrac{16}{9}\right)+\ln\left(\dfrac34\right)$
  \task $\ln\left(e^{3}\right)+\ln\left(e^{-2}\right)-\ln e$
\end{tasks}
\end{exercice}

\begin{exercice}[1]{Puissances et racines}
Calculer les valeurs exactes suivantes.
\begin{tasks}(3)
  \task $\ln\sqrt e$
  \task $\ln\left(e^{5}\right)$
  \task $\ln\left(\dfrac{1}{e^{2}}\right)$
  \task $\ln\sqrt[3]{e}$, où ${\sqrt[3]{e}=e^{1/3}}$
  \task $\ln\left(e^{2x-1}\right)$ pour $x=3$
\end{tasks}
\end{exercice}

\rubrique{Entraînement}

\begin{exercice}[2]{Dérivées}
Calculer la dérivée de chacune des fonctions suivantes, en précisant l'intervalle de dérivabilité.
\begin{tasks}(2)
  \task $f(x)=\ln(3x+1)$
  \task $g(x)=\ln\left(x^{2}-4\right)$, sur $\intOO{2}{+\infty}$
  \task $h(x)=x^{2}\ln x$
  \task $k(x)=\dfrac{\ln x}{x}$
  \task $m(x)=\ln\sqrt{x^{2}+1}$
  \task $p(x)=\ln\left(x^{2}+4\right)$
  \task $q(x)=\ln\left(3x^{2}+2x-1\right)$
\end{tasks}
Étudier ensuite les variations de $p$ sur $\R$.
\end{exercice}

\begin{exercice}[2]{Équations}
Résoudre dans $\R$ les équations suivantes, après avoir déterminé l'ensemble de résolution.
\begin{tasks}(2)
  \task $\ln(x+2)=\ln(3x-1)$
  \task $\ln x=2$
  \task $\ln\left(x^{2}-4\right)-\ln(x-2)=1$
  \task $\ln(x+1)+\ln(x-1)=\ln8$
  \task $2\ln x=\ln\left(x^{2}+3\right)$
\end{tasks}
\end{exercice}

\begin{exercice}[2]{Inéquations}
Résoudre dans $\R$ les inéquations suivantes.
\begin{tasks}(2)
  \task $\ln(x+1)<\ln(2x)$
  \task $\ln x\geq-1$
  \task $\ln(2x+3)\leq0$
  \task $\ln\left(x^{2}-1\right)>\ln3$
  \task $\ln x+\ln(x-2)\leq\ln8$
\end{tasks}
\end{exercice}

\begin{exercice}[2]{Limites et croissances comparées}
Calculer les limites suivantes.
\begin{tasks}(3)
  \task $\displaystyle\lim_{x\to0^{+}}x\ln x$
  \task $\displaystyle\lim_{x\to+\infty}\frac{\ln x}{x}$
  \task $\displaystyle\lim_{x\to0}\frac{\ln(1+x)}{x}$
  \task $\displaystyle\lim_{x\to+\infty}\frac{x}{\ln x}$
  \task $\displaystyle\lim_{x\to1}\frac{\ln x}{x-1}$
  \task $\displaystyle\lim_{x\to+\infty}\frac{\ln\left(x^{2}\right)}{x}$
\end{tasks}
\end{exercice}

\begin{exercice}[2]{Études de fonctions}
Pour chacune des fonctions suivantes, déterminer l'ensemble de définition, la dérivée, le sens de variation
et les limites aux bornes de l'ensemble de définition.
\begin{tasks}(3)
  \task $f(x)=\ln(2x-1)$
  \task $g(x)=x\ln x$
  \task $h(x)=\ln\left(x^{2}-3x+2\right)$
  \task $k(x)=\dfrac{\ln x}{x^{2}}$
  \task $m(x)=\ln\left(\dfrac{x+1}{x-1}\right)$
\end{tasks}
\end{exercice}

\rubrique{Problèmes}

\begin{exercice}[2]{Modélisation}
\begin{questions}
  \item Une population de bactéries suit la loi $P(t)=P_0e^{kt}$, avec $P_0=100$ et $t$ en heures. Elle double en
    $3$ heures. Déterminer $k$, puis le temps nécessaire pour atteindre $10\,000$ bactéries.
  \item Un capital de $1\,000$\,€ est placé à intérêts composés au taux annuel de 5\,\%. Au bout de combien d'années
    aura-t-il doublé ?
  \item La température d'un objet suit la loi $T(t)=20+80e^{-kt}$ ($t$ en minutes). Sachant que $T(10)=60$, déterminer $k$.
  \item Une substance radioactive se désintègre selon $Q(t)=Q_0e^{-\lambda t}$ ($t$ en années). Sa demi-vie est de
    $30$ ans. Calculer $\lambda$, puis la proportion de substance restante au bout de $100$ ans.
  \item Dans un modèle d'intensité sonore, le niveau $L$ est lié à l'intensité $I$ par $L=10\ln\left(\dfrac{I}{I_0}\right)$.
    Exprimer $I$ en fonction de $I_0$ lorsque $L=80$.
\end{questions}
\end{exercice}

\begin{exercice}[3]{Synthèse 1}
\begin{questions}
  \item Étudier la fonction $f$ définie par $f(x)=\ln\left(x^{2}-4x+3\right)$ : ensemble de définition, variations, limites.
  \item Résoudre le système $\ln x+\ln y=\ln12$ et $\ln x-\ln y=\ln3$, d'inconnues $x>0$ et $y>0$.
  \item Montrer que, pour tout $x>-1$, $\ln(1+x)\leq x$.
  \item Calculer $\displaystyle\lim_{x\to+\infty}\frac{\ln\left(x^{2}+1\right)}{x}$.
  \item Déterminer les réels $a$ pour lesquels l'équation $\ln x=ax$ admet exactement une solution. On pourra étudier
    la fonction $x\mapsto\dfrac{\ln x}{x}$.
\end{questions}
\end{exercice}

\begin{exercice}[3]{Synthèse 2}
On considère la fonction $f$ définie sur $\intOO{0}{+\infty}$ par $f(x)=\ln x-x+2$.
\begin{questions}
  \item Calculer $f'(x)$ et étudier son signe.
  \item Calculer les limites de $f$ en $0$ et en $+\infty$, puis dresser son tableau de variations.
  \item Déterminer le nombre de solutions de l'équation $\ln x=x-2$, et en donner des valeurs approchées.
  \item Résoudre l'inéquation $\ln x\geq x-2$.
\end{questions}
\end{exercice}

\begin{exercice}[3]{Une composée avec le logarithme}
On considère la fonction $f$ définie par $f(x)=\ln\left(\dfrac{x^{2}+1}{x^{2}-4}\right)$.
\begin{questions}
  \item Justifier que $f$ est définie sur $D=\intOO{-\infty}{-2}\cup\intOO{2}{+\infty}$, et que $f$ est paire.
  \item Montrer que, pour tout $x\in D$, $f'(x)=\dfrac{-10x}{\left(x^{2}+1\right)\left(x^{2}-4\right)}$, puis étudier les variations de $f$.
  \item Calculer les limites de $f$ aux bornes de $D$. Préciser les asymptotes à la courbe.
  \item Montrer que $f(x)>0$ pour tout $x\in D$.
  \item On admet que $f''(x)=\dfrac{10\left(3x^{4}-3x^{2}+4\right)}{\left(x^{2}+1\right)^{2}\left(x^{2}-4\right)^{2}}$. Étudier la convexité de $f$.
\end{questions}
\end{exercice}

\begin{exercice}[3]{Fonction auxiliaire et asymptote oblique}
\partieExo{A}{Une fonction auxiliaire}
On considère la fonction $g$ définie sur $\intOO{0}{+\infty}$ par $g(x)=x^{2}-1+\ln x$.
\begin{questions}
  \item Étudier les variations de $g$.
  \item Calculer $g(1)$, puis en déduire le signe de $g(x)$ selon les valeurs de $x$.
\end{questions}
\partieExo{B}{Étude de la fonction $f$}
On considère la fonction $f$ définie sur $\intOO{0}{+\infty}$ par $f(x)=x-1-\dfrac{\ln x}{x}$, de courbe $\Cf$.
\begin{questions}[resume]
  \item Calculer les limites de $f$ en $0$ et en $+\infty$.
  \item Montrer que, pour tout $x>0$, $f'(x)=\dfrac{g(x)}{x^{2}}$. En déduire le tableau de variations de $f$.
  \item Montrer que la droite $\Delta$ d'équation $y=x-1$ est asymptote à $\Cf$ en $+\infty$, et étudier leur position relative.
  \item Tracer $\Delta$ et $\Cf$.
\end{questions}
\end{exercice}

\begin{exercice}[3]{Injections répétées d'un médicament}
\partieExo{A}{Une seule injection}
On injecte $20$~mg d'un médicament à un patient. Chaque heure, l'organisme élimine 30\,\% de la quantité présente.
On note $u_n$ la quantité (en mg) présente au bout de $n$ heures.
\begin{questions}
  \item Exprimer $u_n$ en fonction de $n$.
  \item À l'aide du logarithme, déterminer au bout de combien d'heures il reste moins de $1$~mg de médicament.
\end{questions}
\partieExo{B}{Injections toutes les heures}
On injecte maintenant $5$~mg à l'instant $0$, puis, au bout de chaque heure, 30\,\% du médicament est éliminé et l'on
réinjecte $5$~mg. On note $v_n$ la quantité présente juste après la $n$-ième réinjection, avec $v_0=5$.
\begin{questions}[resume]
  \item Justifier que $v_{n+1}=0{,}7\,v_n+5$, puis montrer que la suite $w_n=v_n-\dfrac{50}{3}$ est géométrique.
  \item En déduire que $v_n=\dfrac{50}{3}-\dfrac{35}{3}\times0{,}7^{n}$, puis la limite de $(v_n)$. Interpréter.
  \item Le traitement est efficace lorsque la quantité présente est au moins égale à $15$~mg. Déterminer, à l'aide du
    logarithme, le nombre minimal de réinjections nécessaires.
\end{questions}
\end{exercice}

\begin{exercice}[3]{La série harmonique diverge}
Pour tout entier $n\geq1$, on pose $H_n=1+\dfrac12+\dfrac13+\dots+\dfrac1n$.
\partieExo{A}{Une inégalité}
\begin{questions}
  \item Montrer que, pour tout $x>0$, $\ln x\leq x-1$. On pourra étudier la fonction $x\mapsto x-1-\ln x$.
  \item En déduire que, pour tout entier $k\geq1$, $\ln\left(1+\dfrac1k\right)\leq\dfrac1k$.
\end{questions}
\partieExo{B}{Divergence de $(H_n)$}
\begin{questions}[resume]
  \item Vérifier que $\ln\left(1+\dfrac1k\right)=\ln(k+1)-\ln k$, puis calculer $\displaystyle\sum_{k=1}^{n}\ln\left(1+\frac1k\right)$.
  \item En déduire que, pour tout $n\geq1$, $H_n\geq\ln(n+1)$, puis la limite de la suite $(H_n)$.
  \item Déterminer un entier $n$ à partir duquel on est certain que $H_n>10$.
  \item Que renvoie la fonction Python suivante ? Que constate-t-on sur la vitesse de croissance de $(H_n)$ ?
\begin{python}
def seuil():
    n = 0
    h = 0
    while h <= 5:
        n = n + 1
        h = h + 1/n
    return n
\end{python}
\end{questions}
\end{exercice}
